\section{Reproduction map and environment}\label{app:repro}
All code is in \texttt{code/} (library \texttt{gate/}); see \texttt{README.md} for commands. Python 3.11.2, numpy 2.4.6, scipy 1.17.1, networkx 3.6.1, pandas 3.1.0, matplotlib 3.11.2, scikit-learn 1.9.1. Graphs: SNAP ego-Facebook (\texttt{facebook\_combined.txt.gz}, sha256 \texttt{125e84db\ldots ad035d7}, 4039 nodes, 88{,}234 edges; \citealt{mcauley2012}) and Network Repository \texttt{socfb-Stanford3} (zip sha256 \texttt{b1221e41\ldots 5e8720d}; 11{,}586 nodes, 568{,}309 edges, CC BY-SA; \citealt{rossi2015,traud2012}); both downloaded on 2026-11-01 and checked against the hashes. The largest connected component is used (both graphs are connected). Population seed $0$; replicate $r$ of a configuration uses seed $[2026,\mathrm{crc32}(\mathrm{config}),r]$. Machine: 4 usable CPU cores. Wall-clock: Exp.~1 36 s, Exp.~2 25 min, Exp.~3 22 min, Exp.~4 4.3 min, Exp.~5 13 min, Exp.~6 20 min, Exp.~7 9 s, Exp.~8 13 s, Exp.~9 7 min, Exp.~10 13 min, Exp.~11 23 s.

\begin{table}[h]
\caption{Experiment $\to$ script $\to$ artifact (all scripts in \texttt{code/}; run \texttt{python experiments.py <exp>} first).}\label{tab:repro}
\centering\scriptsize
\setlength{\tabcolsep}{2pt}
\begin{tabular}{lll}
\toprule
Exp. (name) & Script & Artifacts\\
\midrule
1 (\texttt{e1\_var}) & \texttt{fig\_e1.py} & Fig.~\ref{fig:e1}\\
2 (\texttt{e2\_main}) & \texttt{fig\_e2.py}, \texttt{tab\_e2.py} & Table~\ref{tab:e2}, Fig.~\ref{fig:e2app}\\
3 (\texttt{e3\_design}) & \texttt{fig\_e3.py} & Fig.~\ref{fig:e3}, Table~\ref{tab:e3}\\
4 (\texttt{e4\_scale}) & \texttt{fig\_e4.py} & Fig.~\ref{fig:e4}\\
5 (\texttt{e5\_dosenoise}) & \texttt{fig\_e5.py} & Fig.~\ref{fig:e5}\\
6 (\texttt{e6\_degree}) & \texttt{fig\_e6.py} & Fig.~\ref{fig:e6}\\
7 (\texttt{e7\_infer}) & \texttt{tab\_e7.py} & Tables~\ref{tab:e7}, \ref{tab:e7full}\\
8 (\texttt{e8\_misspec\_graph}) & \texttt{fig\_e8.py} & Fig.~\ref{fig:e8}\\
9 (\texttt{e9\_prob}) & \texttt{fig\_e9.py} & Fig.~\ref{fig:e9}\\
10 (\texttt{e10\_clusterres}) & \texttt{fig\_e10.py}, \texttt{tab\_app.py} & Table~\ref{tab:e10}\\
11 (\texttt{e11\_homophily}) & \texttt{tab\_app.py} & Table~\ref{tab:e11}\\
1--11 & \texttt{fig\_diag.py} & Fig.~\ref{fig:diag}, R2/R9 numbers\\
6, 9 & \texttt{stats\_theory.py}, \texttt{tab\_theory.py} & Table~\ref{tab:theory}, R3 numbers\\
\bottomrule
\end{tabular}
\end{table}

\section{Implementation details}
\textbf{Estimators.} OLS uses a tiny ridge ($10^{-8}n$ on the diagonal) for numerical stability only. Feature sets: x $=(z,x)$; c $=(z,\mathrm{count}/\bar d)$; xd $=(z,x,\mathrm{zlogdeg})$; cd $=(z,\mathrm{count}/\bar d,\mathrm{zlogdeg})$; x2 $=(z,x,Wx,\mathrm{zlogdeg})$; quad $=(z,x,x^2,zx,\mathrm{zlogdeg})$; rich $=(z,x,x^2,x^3,zx,\mathrm{zlogdeg},z\,\mathrm{zlogdeg},x\,\mathrm{zlogdeg},Wx)$; oracle uses the true exposure basis (plus zlogdeg; for \texttt{hetero} it also uses the latent $u$, for \texttt{peer} the exact $(I-\rho W)^{-1}z$). Note that OLS-x omits the degree covariate, while the baseline contains a degree term. All regressions predict at $z{=}\mathbf 1$ and $z{=}\mathbf 0$ by recomputing the features on the all-treated and all-control vectors. GBDT: scikit-learn \texttt{HistGradientBoostingRegressor} (150 iterations, learning rate 0.1, max depth 4, min leaf 20) on $(z,x,\mathrm{count}/\bar d,Wx,\mathrm{zlogdeg})$, not tuned. H\'ajek: exposure indicators $\{z_i{=}1,x_i\ge0.75\}$ and $\{z_i{=}0,x_i\le0.25\}$, probabilities by 400 Monte-Carlo design draws (floored at $1/401$), H\'ajek-normalised; undefined when either group is empty, counted as failure (and excluded from RMSE; failure rates are in the summary CSVs).\\
\textbf{Mechanism parameters.} $\tau_d{=}1$, $\gamma{=}1$, \texttt{sat}: $\kappa{=}4$; \texttt{thresh}: centre $0.7$, slope $10$; \texttt{peer}: $\rho{=}0.5$; \texttt{hetero}: $(1+0.5u_i)$ with $u=0.6\,\mathrm{zlogdeg}+0.8\,\xi$ standardised; baseline $a_i=1+0.5\,\mathrm{zlogdeg}_i+s_b\,b_{c(i)}+0.7\nu_i$, $s_b{=}0.5$ ($c(i)$: Louvain community at resolution 1). The GATE is computed exactly; it equals $2.000$ for all mechanisms except \texttt{hetero}/\texttt{peer} where it is the computed value (also $\approx2$). The ``noiseless refit'' gives the assignment-induced part $S$.\\
\textbf{Clusters.} Louvain \citep{blondel2008} via networkx with seed 0; resolution per graph for the main grid: ER 3 ($K{=}110$, cut 0.82), BA 3 (111, 0.82), WS 3 (49, 0.16), SBM 1 (30, 0.10), egoFB 3 (27, 0.07), Stanford3 3 (57, 0.64), chosen once after inspecting $K$ and cut and before looking at any outcome. Resolution 1 on ER/BA gives $K\approx10$--12; we did not tune further. Cluster assignment treats $\mathrm{round}(pK)$ clusters.\\
\textbf{Graph corruption (Exp.~8).} A fraction $r$ of edges is deleted uniformly and the same number of uniformly random pairs is added; features are computed on this graph, outcomes on the true graph.\\
\textbf{Inference (Exp.~7).} With $e=y-X\hat\beta$: HC0 $=\sqrt{\sum w_i^2e_i^2}$, HC3 uses $e_i/(1-h_{ii})$, CR $=\sqrt{\sum_g(\sum_{i\in g}w_ie_i)^2}$ over design clusters (Louvain, same as the cluster design; for Bernoulli designs the same partition is used only for the SE), i.i.d.\ $=\hat\sigma\|w\|_2$; intervals $\hat\tau\pm1.96\,\mathrm{SE}$ vs.\ the exact GATE.

\section{Additional results}
\begin{figure*}[t]
\centering
\includegraphics[width=0.9\textwidth]{../figs/fig_e2_app.pdf}
\caption{Relative RMSE of the other estimators in the main grid (Exp.~2): rows are graph and design (B: Bernoulli, C: cluster), columns mechanisms. 200 replicates per cell.}\label{fig:e2app}
\end{figure*}
\begin{figure*}[t]
\centering
\includegraphics[width=0.9\textwidth]{../figs/fig_e2_main.pdf}
\caption{Heatmap version of the main grid (DM, OLS-x, GBDT); relative RMSE.}\label{fig:e2main}
\end{figure*}
\begin{figure}[t]
\centering
\includegraphics[width=\columnwidth]{../figs/fig_e4.pdf}
\caption{Scaling with $n$ (SBM, mean degree 16, 100 nodes per block, 100 replicates). Solid: Bernoulli, dashed: cluster. Top: RMSE; bottom: $|$bias$|$. Dotted: slope $-1/2$. The GBDT bias and the OLS-x bias for \texttt{sat}/\texttt{thresh} under Bernoulli do not vanish with $n$.}\label{fig:e4}
\end{figure}
\begin{figure}[t]
\centering
\includegraphics[width=\columnwidth]{../figs/fig_e5.pdf}
\caption{Interference strength $\gamma$ and noise $\sigma$ (SBM, 200 replicates). Solid: Bernoulli; dashed: cluster. Rows: mechanism.}\label{fig:e5}
\end{figure}

\begin{table}[t]
\caption{Prop.~\ref{prop:bern} check, SBM graph, Bernoulli($p$): empirical bias of the noiseless OLS-x fit ($\pm$ Monte-Carlo s.e.) vs.\ the finite-degree prediction using the realised degree sequence (Exp.~9, 200 replicates).}\label{tab:theory}
\centering\scriptsize
\begin{tabular}{lrrrrr}
\toprule
Mech. & $p$ & empirical & s.e. & predicted & diff \\
\midrule
linfrac & 0.05 & -0.007 & 0.055 & +0.000 & -0.007 \\
linfrac & 0.10 & +0.023 & 0.040 & -0.000 & +0.023 \\
linfrac & 0.20 & +0.016 & 0.031 & -0.000 & +0.016 \\
linfrac & 0.30 & -0.010 & 0.026 & +0.000 & -0.010 \\
linfrac & 0.50 & +0.029 & 0.024 & -0.000 & +0.029 \\
sat & 0.05 & +2.049 & 0.059 & +2.016 & +0.033 \\
sat & 0.10 & +1.566 & 0.041 & +1.556 & +0.009 \\
sat & 0.20 & +0.858 & 0.030 & +0.826 & +0.032 \\
sat & 0.30 & +0.291 & 0.026 & +0.293 & -0.002 \\
sat & 0.50 & -0.365 & 0.023 & -0.369 & +0.004 \\
thresh & 0.05 & -0.901 & 0.055 & -0.973 & +0.072 \\
thresh & 0.10 & -0.960 & 0.041 & -0.950 & -0.011 \\
thresh & 0.20 & -0.832 & 0.030 & -0.852 & +0.020 \\
thresh & 0.30 & -0.630 & 0.024 & -0.642 & +0.012 \\
thresh & 0.50 & +0.248 & 0.023 & +0.231 & +0.018 \\
\bottomrule
\end{tabular}

\end{table}

\begin{table*}[t]
\caption{Cluster granularity (Exp.~10): RMSE for \texttt{sat} (left four) and \texttt{thresh} (right four). Res.: Louvain resolution; $K$ clusters; cut: fraction of edges between clusters; cluster design with half of the clusters treated; 200 replicates.}\label{tab:e10}
\centering\scriptsize
\begin{tabular}{llrr|rrrr|rrrr}
\toprule
Graph & Res. & $K$ & cut & DM & OLS-x & Orac. & GBDT & DM & OLS-x & Orac. & GBDT \\
\midrule
egoFB & 1 & 16 & 0.04 & 0.34 & 0.34 & 0.37 & 0.97 & 0.34 & 0.35 & 0.36 & 0.95 \\
egoFB & 3 & 27 & 0.07 & 0.36 & 0.41 & 0.41 & 0.99 & 0.33 & 0.39 & 0.37 & 1.06 \\
egoFB & 10 & 238 & 0.53 & 0.40 & 0.26 & 0.29 & 1.13 & 0.29 & 0.26 & 0.23 & 1.12 \\
egoFB & 30 & 672 & 0.73 & 0.56 & 0.31 & 0.29 & 1.06 & 0.45 & 0.30 & 0.22 & 1.12 \\
\midrule
Stanford3 & 1 & 10 & 0.46 & 0.77 & 0.60 & 0.66 & 0.61 & 0.68 & 0.60 & 0.61 & 0.56 \\
Stanford3 & 3 & 57 & 0.64 & 0.66 & 0.49 & 0.57 & 0.54 & 0.53 & 0.64 & 0.53 & 0.54 \\
Stanford3 & 10 & 197 & 0.74 & 0.70 & 0.41 & 0.52 & 0.47 & 0.58 & 0.58 & 0.44 & 0.46 \\
\midrule
WS & 1 & 27 & 0.13 & 0.32 & 0.24 & 0.23 & 0.62 & 0.22 & 0.22 & 0.21 & 0.65 \\
WS & 3 & 49 & 0.16 & 0.28 & 0.19 & 0.18 & 0.67 & 0.21 & 0.20 & 0.18 & 0.66 \\
WS & 10 & 89 & 0.21 & 0.33 & 0.20 & 0.18 & 0.65 & 0.20 & 0.18 & 0.15 & 0.57 \\
WS & 30 & 156 & 0.30 & 0.41 & 0.17 & 0.15 & 0.54 & 0.25 & 0.17 & 0.12 & 0.60 \\
\midrule
SBM(0.3) & 1 & 30 & 0.31 & 0.50 & 0.23 & 0.22 & 0.42 & 0.33 & 0.24 & 0.23 & 0.41 \\
SBM(0.3) & 3 & 30 & 0.31 & 0.50 & 0.23 & 0.22 & 0.41 & 0.30 & 0.23 & 0.21 & 0.41 \\
SBM(0.3) & 30 & 376 & 0.78 & 0.86 & 0.49 & 0.45 & 0.76 & 0.73 & 0.48 & 0.28 & 0.66 \\
\bottomrule
\end{tabular}

\end{table*}

\begin{table}[t]
\caption{Community-level baseline effect sd $s_b\in\{0,0.5,1,2\}$ (Exp.~11, \texttt{sat}): RMSE (GATE $=2$), 300 replicates.}\label{tab:e11}
\centering\scriptsize
\begin{tabular}{lllrrrr}
\toprule
Graph & Des. & Est. & 0 & 0.5 & 1 & 2 \\
\midrule
SBM & Bern & DM & 1.00 & 1.00 & 1.00 & 1.01 \\
SBM & Bern & OLS-x & 0.41 & 0.51 & 0.74 & 1.37 \\
SBM & Bern & OLS-oracle & 0.25 & 0.52 & 0.89 & 1.77 \\
SBM & Clust & DM & 0.18 & 0.28 & 0.47 & 0.87 \\
SBM & Clust & OLS-x & 0.06 & 0.21 & 0.41 & 0.84 \\
SBM & Clust & OLS-oracle & 0.05 & 0.20 & 0.41 & 0.85 \\
\midrule
WS & Bern & DM & 1.00 & 1.00 & 1.00 & 1.00 \\
WS & Bern & OLS-x & 0.42 & 0.55 & 0.73 & 1.41 \\
WS & Bern & OLS-oracle & 0.27 & 0.53 & 0.94 & 1.92 \\
WS & Clust & DM & 0.24 & 0.28 & 0.38 & 0.69 \\
WS & Clust & OLS-x & 0.07 & 0.19 & 0.37 & 0.78 \\
WS & Clust & OLS-oracle & 0.05 & 0.18 & 0.35 & 0.75 \\
\midrule
egoFB & Bern & DM & 1.01 & 1.00 & 1.00 & 1.00 \\
egoFB & Bern & OLS-x & 0.66 & 0.62 & 0.90 & 1.63 \\
egoFB & Bern & OLS-oracle & 0.15 & 0.45 & 0.85 & 1.61 \\
egoFB & Clust & DM & 0.22 & 0.33 & 0.64 & 1.37 \\
egoFB & Clust & OLS-x & 0.22 & 0.37 & 0.75 & 1.63 \\
egoFB & Clust & OLS-oracle & 0.04 & 0.37 & 0.77 & 1.60 \\
\bottomrule
\end{tabular}

\end{table}

\begin{table*}[t]
\caption{Full coverage table (Exp.~7; columns as in Table~\ref{tab:e7}).}\label{tab:e7full}
\centering\scriptsize
\begin{tabular}{lllrrrrrrr}
\toprule
Graph & Mech. & Des. & bias & RMSE & iid & HC3 & CR & $\overline{\rm SE}_{iid}$/SD & $\overline{\rm SE}_{CR}$/SD \\
\midrule
SBM & hetero & Bern & -0.05 & 0.37 & 0.74 & 0.76 & 0.94 & 0.61 & 1.01 \\
SBM & hetero & Clust & -0.02 & 0.22 & 0.48 & 0.48 & 0.93 & 0.34 & 0.97 \\
SBM & linfrac & Bern & +0.01 & 0.36 & 0.76 & 0.76 & 0.93 & 0.59 & 0.97 \\
SBM & linfrac & Clust & +0.00 & 0.22 & 0.48 & 0.48 & 0.92 & 0.32 & 0.93 \\
SBM & peer & Bern & -0.41 & 0.76 & 0.47 & 0.47 & 0.87 & 0.39 & 0.93 \\
SBM & peer & Clust & -0.10 & 0.40 & 0.34 & 0.34 & 0.94 & 0.22 & 0.99 \\
SBM & sat & Bern & -0.35 & 0.50 & 0.55 & 0.56 & 0.82 & 0.58 & 0.98 \\
SBM & sat & Clust & -0.02 & 0.23 & 0.45 & 0.46 & 0.90 & 0.31 & 0.91 \\
ER & hetero & Bern & -0.07 & 0.28 & 0.92 & 0.93 & 0.88 & 0.88 & 0.89 \\
ER & hetero & Clust & -0.07 & 0.32 & 0.80 & 0.80 & 0.82 & 0.67 & 0.73 \\
ER & linfrac & Bern & +0.01 & 0.27 & 0.91 & 0.91 & 0.89 & 0.84 & 0.84 \\
ER & linfrac & Clust & +0.01 & 0.32 & 0.79 & 0.79 & 0.82 & 0.64 & 0.69 \\
ER & peer & Bern & -0.47 & 0.57 & 0.49 & 0.50 & 0.53 & 0.75 & 0.78 \\
ER & peer & Clust & -0.41 & 0.57 & 0.48 & 0.48 & 0.54 & 0.54 & 0.60 \\
ER & sat & Bern & -0.37 & 0.47 & 0.61 & 0.63 & 0.59 & 0.80 & 0.81 \\
ER & sat & Clust & -0.35 & 0.46 & 0.56 & 0.57 & 0.60 & 0.68 & 0.74 \\
\bottomrule
\end{tabular}

\end{table*}

\section{Process log, access statement and checklist}
\textbf{Process notes (honest account).} (i) A first launch used 38 worker processes although only 4 cores were available (cgroup limit); it was killed and restarted with 4 workers (no results from it were used). (ii) Exp.~6 first crashed for $\bar d{=}128$ on the 30-block SBM (edge probability above 1); it was re-run with 15 blocks of 200 nodes, and the corrected configuration is the one reported. (iii) Exp.~10 and 11 were added after the first results to address cluster granularity and homophily. (iv) An exploratory smoke test (\texttt{e0\_smoke}, 20 replicates) preceded the main runs and was not used. (v) Design choices that were made after seeing smoke-test output: dropping the count-feature estimator \texttt{c} from the headline comparison (it is confounded with degree in our baseline), adding the community baseline-effect ablation (Exp.~1, 11). (vi) No hyper-parameter was tuned on outcomes. Nothing in the paper comes from a run that is not reproducible with the shipped code, and every reported number is regenerated by the scripts above.

\textbf{Restricted content.} The task excluded one published paper and its derivatives. We did not search for it; our literature searches were by the exact titles of well-known papers (listed in the references) and the starter list. We did not knowingly encounter any excluded material.

\textbf{Checklist.} Claims match evidence: yes (Sec.~\ref{sec:res}; limits in Sec.~\ref{sec:disc}). Theory assumptions and proofs: Props.~\ref{prop:decomp}--\ref{prop:bern} state assumptions; Prop.~\ref{prop:bern} is a sketch. Code/data: provided with hashes and commands. Statistical significance: bootstrap CIs and Monte-Carlo s.e. (replicates 100--500; one fixed population per configuration). Compute: 4 CPU cores, $\approx1.9$ h of simulation wall-clock in total. Licences: both graphs are public (CC BY-SA / cited). No human subjects. Impact: see Impact Statement.
